\documentclass{article}

\title{Ch. 31 CW (Cont'd)}
\author{Principles of Macroeconomics}
\date{\today}

\begin{document}

\maketitle

\section*{Instructions}
Please answer the following problems on a separate sheet of paper.

\section*{Problems}

\begin{enumerate}

\item It is sometimes suggested that the Fed should try to achieve zero inflation. If we assume that velocity is constant, does this zero-inflation goal require that the rate of money growth equal zero? If yes, explain why. If no, explain what the rate of money growth should equal.

\item Suppose that a country's inflation rate increases sharply. What happens to the inflation tax on the holders of money? Why is wealth held in savings accounts \textbf{not} subject to a change in the inflation tax? Can you think of any way in which holders of savings accounts are hurt by the increase in inflation?

\item Assuming a tax rate of 40 percent, compute the before-tax real interest rate and the after-tax real interest rate for each of the following cases.

	\begin{enumerate}
	
	\item The nominal interest rate is 10 percent, and the inflation rate is 5 percent.
	
	\item The nominal interest rate is 6 percent, and the inflation rate is 2 percent.
	
	\item The nominal interest rate is 4 percent, and the inflation rate is 1 percent.
	
	\end{enumerate}
	
\item Recall that money serves three functions in the economy. What are those functions? How does inflation affect the ability of money to serve each of these functions?

\newpage

\item Suppose that people expect inflation to be 3 percent but that, in fact, prices rise by 5 percent. Describe how this unexpectedly high inflation would help or hurt the following:

	\begin{enumerate}
	
	\item the government
	
	\item a homeowner with a fixed-rate mortgage
	
	\item a union worker in the second year of a labor contract
	
	\item a college that has invested some of its endowment in government bonds
	
	\end{enumerate}
	
\item Explain whether the following statements are true false or uncertain.

	\begin{enumerate}

	\item ``Inflation hurts borrowers and helps lenders because borrowers must pay a higher rate of interest."
	
	\item ``If prices change in a way that leaves the overall price level unchanged, then no one is made better or worse off."
	
	\item ``Inflation does not reduce the purchasing power of most workers."
	
	\end{enumerate}

\end{enumerate}

\end{document}